% ===== PRÉAMBULE CANONIQUE — à coller VERBATIM en tête de CHAQUE .tex =====
\documentclass[11pt,a4paper]{article}
\usepackage[utf8]{inputenc}\usepackage[T1]{fontenc}\usepackage{lmodern}
\usepackage[french]{babel}
\usepackage{amsmath,amssymb,mathtools}\usepackage{icomma}
\usepackage{siunitx}\sisetup{locale=FR}  % \qty{}{}, \si{}, \ang{}
\usepackage{xcolor}\usepackage{colortbl}
\usepackage{tikz}\usepackage{pgfplots}\pgfplotsset{compat=1.18}
\usepackage[siunitx,europeanresistors]{circuitikz} % circuits (élec.)
\usepackage[most]{tcolorbox}\usepackage{enumitem}\usepackage{array,booktabs}
\usepackage[a4paper,margin=2.2cm]{geometry}
\usetikzlibrary{arrows.meta,calc,angles,quotes,positioning,patterns,%
  backgrounds,shapes.geometric,decorations.pathmorphing,decorations.markings,intersections}
\definecolor{primary}{RGB}{222,49,99}\definecolor{primarydark}{RGB}{150,22,55}
\definecolor{defcol}{RGB}{0,114,178}\definecolor{methcol}{RGB}{52,92,148}
\definecolor{excol}{RGB}{0,135,90}\definecolor{attncol}{RGB}{200,40,55}
\tcbset{boxbase/.style={enhanced,breakable,boxrule=0.9pt,arc=2.5pt,
  left=3mm,right=3mm,top=2.3mm,bottom=2.3mm,fonttitle=\bfseries,coltitle=white}}
% --- boîtes TUTORIEL (noms EXACTS, ne pas renommer) ---
\newtcolorbox{cle}[1][]{boxbase,colback=defcol!6,colframe=defcol,
  title={\textsc{À retenir}\ifx\relax#1\relax\relax\else\textmd{ : \itshape #1}\fi}}
\newtcolorbox{methode}[1][]{boxbase,colback=methcol!7,colframe=methcol,sharp corners,
  title={\textsc{Méthode}\ifx\relax#1\relax\relax\else\textmd{ --- \itshape #1}\fi}}
\newtcolorbox{exemplebox}[1][]{boxbase,colback=excol!5,colframe=excol,
  title={\textsc{Exemple}\ifx\relax#1\relax\relax\else\textmd{ : \itshape #1}\fi}}
\newtcolorbox{piege}[1][]{boxbase,colback=attncol!6,colframe=attncol,
  title={\textsc{Piège}\ifx\relax#1\relax\relax\else\textmd{ : \itshape #1}\fi}}
\newtcolorbox{remarque}[1][]{boxbase,colback=black!4,colframe=black!45,
  title={\textsc{Remarque}\ifx\relax#1\relax\relax\else\textmd{ : \itshape #1}\fi}}
% --- boîtes EXERCICE / CORRIGÉ (noms EXACTS, auto-comptées) ---
\newtcolorbox[auto counter]{exo}[1][]{boxbase,colback=primary!4,colframe=primary,
  title={\textsc{Exercice}~\thetcbcounter\ifx\relax#1\relax\relax\else\textmd{ --- \itshape #1}\fi}}
\newtcolorbox[auto counter]{corrige}[1][]{boxbase,colback=excol!4,colframe=excol,
  title={\textsc{Corrigé}~\thetcbcounter\ifx\relax#1\relax\relax\else\textmd{ --- \itshape #1}\fi}}
\newcommand{\vv}[1]{\vec{#1}}\newcommand{\ds}{\displaystyle}
\newcommand{\R}{\mathbb{R}}\newcommand{\N}{\mathbb{N}}\newcommand{\Z}{\mathbb{Z}}
\begin{document}
% THEME_TITLE: Relation de conjugaison et grandissement
\sisetup{per-mode=symbol}

Pour connaître la \emph{position} et la \emph{taille} de l'image \textbf{sans dessiner}, deux formules
suffisent. Toute la difficulté est dans les \textbf{signes} : respecte le tableau et elles marchent
dans tous les cas.

\begin{cle}[les deux formules et leurs conventions]
\[ \boxed{\ \dfrac{1}{f}=\dfrac{1}{u}+\dfrac{1}{v}\ }\qquad\qquad
   \boxed{\ \gamma=\dfrac{i}{o}=-\dfrac{v}{u}\ } \]
\begin{center}\small
\renewcommand{\arraystretch}{1.25}
\begin{tabular}{>{\bfseries}l c c}
\toprule
Grandeur & \cellcolor{excol!12}signe $+$ & \cellcolor{attncol!12}signe $-$ \\
\midrule
$u$ (objet--lentille) & objet réel (cas usuel) & --- \\
$v$ (image--lentille) & image \textbf{réelle} & image \textbf{virtuelle} \\
$f$ (focale) & lentille \textbf{convergente} & lentille \textbf{divergente} \\
$i$ (hauteur image) & image \textbf{droite} & image \textbf{renversée} \\
\bottomrule
\end{tabular}
\end{center}
Ici $o>0$ (objet vers le haut), $u=\overline{OA}$, $v=\overline{OA'}$, $i$ = hauteur de l'image.
\end{cle}

\begin{center}
\begin{tikzpicture}[scale=0.8,>=Stealth,font=\footnotesize]
  \draw[black!35,dash pattern=on 3pt off 2pt] (-6.2,0)--(4.6,0);
  \draw[line width=1.5pt,<->,black!65] (0,-1.9)--(0,1.9);
  \fill[black] (0,0) circle (1.2pt); \node[above right] at (0,0) {$O$};
  \fill[primary] (2,0) circle (1.4pt); \node[above] at (2,0) {$F'$};
  \fill[primary] (-2,0) circle (1.4pt); \node[above] at (-2,0) {$F$};
  \draw[attncol,line width=1.8pt,->] (-5,0)--(-5,1.4); \node[attncol,left] at (-5,0.7) {$o$};
  \draw[excol,line width=1.8pt,->] (3.33,0)--(3.33,-0.93); \node[excol,right] at (3.33,-0.5) {$i$};
  \draw[black!75,line width=0.8pt] (-5,1.4)--(0,0); \draw[defcol,line width=0.8pt,->] (0,0)--(3.9,-1.09);
  \draw[black!75,line width=0.8pt] (-5,1.4)--(0,1.4); \draw[defcol,line width=0.8pt,->] (0,1.4)--(3.9,-1.33);
  % cotes u et v
  \draw[<->,black!55] (-5,-1.55)--(0,-1.55); \node[below,black!70] at (-2.5,-1.55) {$u>0$};
  \draw[<->,black!55] (0,-1.55)--(3.33,-1.55); \node[below,black!70] at (1.7,-1.55) {$v>0$ (réelle)};
\end{tikzpicture}
\end{center}

\begin{methode}[résoudre un problème de conjugaison]
\begin{enumerate}[leftmargin=1.7em,topsep=2pt,itemsep=1pt]
  \item Mettre \textbf{toutes les longueurs dans la même unité} et fixer les signes ($u>0$ ;
        $f>0$ convergente / $f<0$ divergente).
  \item Isoler l'inconnue : $\dfrac{1}{v}=\dfrac{1}{f}-\dfrac{1}{u}$ pour trouver $v$
        (ou $\dfrac{1}{f}=\dfrac{1}{u}+\dfrac{1}{v}$ pour trouver $f$).
  \item Calculer le grandissement $\gamma=-\dfrac{v}{u}$ puis la taille $i=\gamma\,o$.
  \item \textbf{Interpréter les signes} : $v>0$ réelle / $v<0$ virtuelle ; $\gamma<0$ renversée /
        $\gamma>0$ droite ; $|\gamma|>1$ agrandie / $|\gamma|<1$ réduite.
\end{enumerate}
\end{methode}

\begin{exemplebox}[un calcul complet]
Convergente $f=+\qty{20}{\cm}$, objet réel à $u=+\qty{60}{\cm}$, hauteur $o=\qty{2}{\cm}$.
\[ \frac{1}{v}=\frac{1}{20}-\frac{1}{60}=\frac{3-1}{60}=\frac{1}{30}\ \Rightarrow\ v=+\qty{30}{\cm}\ (\text{réelle}). \]
\[ \gamma=-\frac{v}{u}=-\frac{30}{60}=-0{,}5\ \Rightarrow\ i=\gamma\,o=-\qty{1}{\cm}. \]
Image \textbf{réelle}, \textbf{renversée}, \textbf{réduite} de moitié : c'est l'appareil photo.
\end{exemplebox}

\begin{piege}[trois signes qui coûtent des points]
\begin{itemize}[leftmargin=1.5em,topsep=2pt,itemsep=1pt]
  \item \textbf{Même unité} pour $u$, $v$, $f$ : ne mélange pas cm et m dans une même formule.
  \item N'oublie pas le \textbf{« moins »} de $\gamma=-v/u$ : c'est lui qui renverse l'image réelle.
  \item Une image \textbf{virtuelle} donne $v<0$ (et alors $\gamma>0$, image droite) : c'est le cas de
        la loupe et de \emph{toute} lentille divergente.
\end{itemize}
\end{piege}

\end{document}
